\scanpage{20}

To come back to the general case, maybe I should add that a \uline{necessary} condition for $\unsure{\mathscr{J}|\mathcal{J}|\gamma}/\Gamma$ to be algebraic and the Hodge-structures $\mathscr{V}/\Gamma$ over $\mathscr{J}/\Gamma$ algebraic too, is (\uline{granting} Tate's and Hodge's conjectures): for every $g\in G(\mathbb{R})^{\unsure{0|\circ}}$, and every $\Gamma_0$ of finite index in $G(\mathbb{Z})$, the smallest \inserted{alg.} subgroup $H$ of $G$ such that $H_{\mathbb{R}}$ contains $g[i_0(U_{\mathbb{R}})]g^{-1}$ \uline{and} $\unsure{\Gamma|T}_0$ is $G$ itself.

\deleted{I did \unsure{not|\deleted{not}} try to test directly this condition, even in the case of the \uline{usual} Siegel modular space; if it should be false in that case, as we know that $\mathscr{J}/\Gamma$ \uline{is} algebraic and that so are the $\mathscr{V}/\Gamma$ over it, it would follow that either Hodge's or Tate's conjecture is wrong}

\margin{Question: If $G_{\mathbb{R}}$ is semi-simple without compact factors, is $G(\mathbb{Z})$ always Zariski-dense in $G$\unsure{?!|??}}

By the way, did you make out if for the modular varieties $\mathscr{J}/\Gamma$ of your
\[
\left\{\begin{aligned}
&\text{Boulder talk (which I just read),}\\
&\text{the image of }\mathscr{J}/\Gamma\text{ in the usual Siegel}\\
&\text{modular variety is an algebraic variety,}\\
&\text{or at least constructible?}
\end{aligned}\right.
\]

% archive sheet number: 19
